\documentclass[11pt]{amsart}
\usepackage[T1]{fontenc}
\usepackage[utf8]{inputenc}
\usepackage{lmodern}
\usepackage[margin=1.1in]{geometry}
\usepackage{microtype}
\usepackage{amsmath,amssymb,amsthm,mathtools}
\usepackage{enumitem}
\usepackage{hyperref}
\usepackage{xcolor}
\hypersetup{colorlinks=true,linkcolor=blue!50!black,citecolor=blue!50!black,urlcolor=blue!50!black,pdfauthor={The Clankers},pdftitle={Secondary Note on Split-Zero Globalization}}
\numberwithin{equation}{section}

\newtheorem{theorem}{Theorem}[section]
\newtheorem{proposition}[theorem]{Proposition}
\newtheorem{corollary}[theorem]{Corollary}
\newtheorem{lemma}[theorem]{Lemma}
\theoremstyle{definition}
\newtheorem{definition}[theorem]{Definition}
\newtheorem{remark}[theorem]{Remark}

\newcommand{\B}{\mathbf B}
\newcommand{\Spec}{\operatorname{Spec}}
\newcommand{\Hom}{\operatorname{Hom}}
\newcommand{\End}{\operatorname{End}}
\newcommand{\Aut}{\operatorname{Aut}}
\newcommand{\im}{\operatorname{Im}}
\newcommand{\id}{\operatorname{id}}
\newcommand{\cO}{\mathcal O}
\newcommand{\cP}{\mathcal P}
\newcommand{\N}{\mathbf N}
\newcommand{\Z}{\mathbf Z}
\newcommand{\Q}{\mathbf Q}
\newcommand{\F}{\mathbf F}
\newcommand{\C}{\mathbf C}
\newcommand{\rad}{\operatorname{rad}}
\newcommand{\Cl}{\operatorname{Cl}}
\newcommand{\Diag}{\mathsf{Diag}}
\newcommand{\SMod}{\mathbf{SMod}}
\newcommand{\Mod}{\mathbf{Mod}}
\newcommand{\one}{\mathbf 1}

\title[Secondary Note on Split-Zero Globalization]{A Secondary Note on Split-Zero Globalization\\[4pt]\large Semimodule Classification, Multiplicative Reconstruction, and All-Orders Zeta Deformation}
\author{The Clankers}
\date{\today}

\begin{document}
\begin{abstract}
We extend the baseline globalization theory of the split-zero semiring
\[
S=G(R)=R\sqcup\{\tau\},
\qquad
r\oplus s=r+s,
\quad
r\oplus\tau=r,
\quad
\tau\oplus\tau=\tau,
\]
\[
rs=rs,
\qquad
r\tau=\tau r=\tau,
\qquad
\tau^2=\tau,
\]
for a nonzero commutative ring $R$ with identity. Throughout, we assume the results of the companion note: the ring reflection $p_R$, the support character $\chi_R$, the classification of ideals, prime ideals, congruences, and localizations, and the supported-zero projector $\nu_R(x)=0_Rx$.

The main results are the exact classification of $G(R)$-semimodules as join-semilattice-indexed diagrams of $R$-modules; the identification of the cancellative subcategory with $R$-Mod; the Boolean-cubical form of free semimodules; the product theorem $G(A\times B)\cong G(A)\times_{\B}G(B)$; reconstruction from the multiplicative monoid together with the missing additive relations; the infinite-dimensionality of the multiplicative-monoid spectrum in the UFD case; the contracted monoid-algebra decomposition; conservative ideal/class/factorization results; and the full all-orders deformation theory of the shifted quotient-size zeta sheet, including universal jets, cumulants, arithmetic kernels, non-Eulerianity, Faulhaber interpolation, factorial interpolation, and the exact endpoint defect controlled by a two-point Stone model.
\end{abstract}

\maketitle
\tableofcontents

\section{Standing background and notation}

Throughout $R$ denotes a nonzero commutative ring with identity and
\[
S:=G(R)=R\sqcup\{\tau\},
\qquad
e:=0_R\in R\subset S.
\]
We assume the notation and the proved results of the baseline globalization note, namely:
\begin{enumerate}[label=(\roman*)]
    \item the ring reflection
    \[
    p_R:S\to R,
    \qquad
    p_R(\tau)=0_R,
    \quad
    p_R(r)=r;
    \]
    \item the support character
    \[
    \chi_R:S\to\B,
    \qquad
    \chi_R(\tau)=0,
    \quad
    \chi_R(r)=1;
    \]
    \item the classification of ideals, prime ideals, congruences, and localizations away from nilpotent supported elements;
    \item the supported-zero projector
    \[
    \nu_R:S\to S,
    \qquad
    \nu_R(x)=ex,
    \]
    with image $\{\tau,e\}\cong\B$.
\end{enumerate}
No proof from that baseline package is repeated below.

The only external analytic facts used are standard properties of the Hurwitz zeta function and of Dedekind zeta functions; see the references at the end.

\section{Algebraic refinements beyond the baseline package}

\subsection{Algebraic normal form and universal quotients}

\begin{proposition}\label{prop:normal-form}
There is a canonical semiring isomorphism
\[
G(R)\xrightarrow{\sim}D_R:=\{(r,b)\in R\times\B: b=0\Rightarrow r=0_R\}
=\{(0_R,0)\}\cup (R\times\{1\}).
\]
\end{proposition}

\begin{proof}
Define
\[
\iota:G(R)\to D_R,
\qquad
\iota(\tau)=(0_R,0),
\qquad
\iota(r)=(r,1)
\quad(r\in R).
\]
This is bijective. On $D_R$, componentwise addition and multiplication restrict to
\[
(r,1)+(s,1)=(r+s,1),
\qquad
(r,1)+(0,0)=(r,1),
\qquad
(0,0)+(0,0)=(0,0),
\]
\[
(r,1)(s,1)=(rs,1),
\qquad
(r,1)(0,0)=(0,0),
\qquad
(0,0)^2=(0,0),
\]
which are exactly the defining operations of $G(R)$ transported through $\iota$. 
\end{proof}

\begin{theorem}\label{thm:ring-quotient}
For every commutative ring $A$, composition with $p_R$ induces a bijection
\[
\Hom_{\mathrm{Ring}}(R,A)
\xrightarrow{\sim}
\Hom_{\mathrm{SemiRing}}(G(R),A).
\]
In particular, $R$ is the universal ring quotient of $G(R)$.
\end{theorem}

\begin{proof}
Let $f:G(R)\to A$ be a semiring homomorphism. Since $\tau=0_S$, one has
\[
f(\tau)=0_A.
\]
Also,
\[
f(0_R)=f(0_R\oplus 0_R)=f(0_R)+f(0_R),
\]
so subtraction in the ring $A$ gives $f(0_R)=0_A$. Restrict $f$ to $R$ and call the restriction $g$. For $r,s\in R$,
\[
g(r+s)=f(r\oplus s)=f(r)+f(s)=g(r)+g(s),
\]
\[
g(rs)=f(rs)=f(r)f(s)=g(r)g(s),
\qquad
 g(1_R)=f(1_R)=1_A.
\]
Hence $g:R\to A$ is a ring homomorphism. Since $p_R(\tau)=0_R$ and $p_R(r)=r$, one has
\[
(g\circ p_R)(\tau)=g(0_R)=0_A=f(\tau),
\qquad
(g\circ p_R)(r)=g(r)=f(r).
\]
Thus $f=g\circ p_R$. Uniqueness follows from surjectivity of $p_R$. 
\end{proof}

\begin{corollary}\label{cor:no-splitting}
There is no semiring homomorphism $s:R\to G(R)$ such that
\[
p_R\circ s=\id_R.
\]
\end{corollary}

\begin{proof}
If such $s$ existed, then $s(0_R)=\tau$ because $s$ preserves the semiring zero, and $s(1_R)=1_R$ because $p_R\circ s=\id_R$. Also $p_R(s(-1_R))=-1_R\neq0_R$, so $s(-1_R)$ is supported. Hence
\[
\tau=s(0_R)=s(1_R+(-1_R))=s(1_R)\oplus s(-1_R)=1_R\oplus s(-1_R),
\]
but the right-hand side is supported, contradiction. 
\end{proof}

\begin{theorem}\label{thm:boolean-character-unique}
The support character
\[
\chi_R:G(R)\to\B
\]
is the unique semiring homomorphism from $G(R)$ to $\B$.
\end{theorem}

\begin{proof}
Let $u:G(R)\to\B$ be a semiring homomorphism. Then $u(\tau)=0$ and $u(1_R)=1$. Now
\[
u(0_R)=u(1_R\oplus (-1_R))=u(1_R)+u(-1_R)=1+u(-1_R)=1
\]
in $\B$. If $u(r)=0$ for some $r\in R$, then
\[
1=u(0_R)=u(r0_R)=u(r)u(0_R)=0,
\]
contradiction. Hence $u(r)=1$ for every $r\in R$. Therefore $u=\chi_R$. 
\end{proof}

\begin{corollary}\label{cor:no-ring-to-bool}
There is no semiring homomorphism $R\to\B$.
\end{corollary}

\begin{proof}
If $v:R\to\B$ were such a homomorphism, then
\[
0=v(0_R)=v(1_R+(-1_R))=v(1_R)+v(-1_R)=1+v(-1_R)=1,
\]
contradiction. 
\end{proof}

\begin{corollary}\label{cor:no-unital-bool-section}
There is no unit-preserving semiring homomorphism $\B\to G(R)$.
\end{corollary}

\begin{proof}
If $w:\B\to G(R)$ were unit-preserving, then $w(1)=1_R$. Since $1+1=1$ in $\B$,
\[
1_R=w(1)=w(1+1)=w(1)\oplus w(1)=1_R\oplus 1_R=1_R+1_R.
\]
Subtracting $1_R$ inside the supported copy of $R$ yields $1_R=0_R$, contradiction. 
\end{proof}

\begin{theorem}\label{thm:full-faithfulness}
For nonzero commutative rings $R,R'$, restriction induces a bijection
\[
\Hom_{\mathrm{SemiRing}}(G(R),G(R'))
\xrightarrow{\sim}
\Hom_{\mathrm{Ring}}(R,R').
\]
Consequently,
\[
\End_{\mathrm{SemiRing}}(G(R))\cong \End_{\mathrm{Ring}}(R),
\qquad
\Aut_{\mathrm{SemiRing}}(G(R))\cong \Aut_{\mathrm{Ring}}(R).
\]
\end{theorem}

\begin{proof}
Let $f:G(R)\to G(R')$ be a semiring homomorphism. Since $f$ preserves semiring zero, $f(\tau)=\tau$. By Theorem~\ref{thm:boolean-character-unique}, the composite $\chi_{R'}\circ f$ must equal $\chi_R$. Therefore for every $r\in R$,
\[
\chi_{R'}(f(r))=\chi_R(r)=1,
\]
so $f(r)$ is supported, hence belongs to $R'$. Thus $f$ restricts to a map $R\to R'$. For $r,s\in R$,
\[
f(r+s)=f(r\oplus s)=f(r)\oplus f(s)=f(r)+f(s),
\]
\[
f(rs)=f(r)f(s),
\qquad
f(1_R)=1_{R'}.
\]
Hence $f|_R$ is a ring homomorphism.

Conversely, if $g:R\to R'$ is a ring homomorphism, define
\[
G(g)(\tau)=\tau,
\qquad
G(g)(r)=g(r)
\quad(r\in R).
\]
A direct check shows that $G(g)$ is a semiring homomorphism. The two constructions are inverse. 
\end{proof}

\subsection{Product structure and multiplicative shadows}

\begin{theorem}\label{thm:product-fiber}
For nonzero commutative rings $A,B$,
\[
G(A\times B)\cong G(A)\times_{\B}G(B),
\]
where the fiber product is taken with respect to the support characters.
\end{theorem}

\begin{proof}
Define
\[
\Phi:G(A\times B)\to G(A)\times_{\B}G(B)
\]
by
\[
\Phi(\tau)=(\tau,\tau),
\qquad
\Phi((a,b))=(a,b).
\]
This is well defined because every supported element of $G(A\times B)$ has support character $1$ in both components, while $\tau$ has support character $0$ in both components.

Conversely, an element of the fiber product is either $(\tau,\tau)$ or a pair $(x,y)$ with $\chi_A(x)=\chi_B(y)=1$, hence $x\in A$ and $y\in B$. Define
\[
\Psi(\tau,\tau)=\tau,
\qquad
\Psi(x,y)=(x,y)\in A\times B.
\]
These maps are inverse semiring homomorphisms. 
\end{proof}

\begin{theorem}\label{thm:monoid-reconstruction}
Let $M_S=(S,\cdot)$ be the multiplicative monoid of $S$. Let $T(S)$ be the quotient of the free commutative semiring $\mathbb N[M_S]$ by the semiring congruence generated by
\[
[a]+[b]\sim [a+b]
\qquad (a,b\in R),
\]
and
\[
[x]+[\tau]\sim [x]
\qquad (x\in M_S).
\]
Then
\[
T(S)\cong S.
\]
\end{theorem}

\begin{proof}
Define
\[
\Phi:\mathbb N[M_S]\to S,
\qquad
\Phi([x])=x.
\]
The generating relations hold in $S$, because addition of supported elements is ring addition and $x\oplus\tau=x$ for every $x\in M_S$. Hence $\Phi$ factors through a semiring homomorphism
\[
\overline\Phi:T(S)\to S.
\]
Conversely, define
\[
\Psi:S\to T(S),
\qquad
\Psi(x)=[x].
\]
This is multiplicative by construction and additive because
\[
\Psi(a\oplus b)=\Psi(a+b)=[a+b]=[a]+[b]
\qquad(a,b\in R),
\]
and
\[
\Psi(x\oplus\tau)=\Psi(x)=[x]=[x]+[\tau]
\qquad(x\in M_S).
\]
The maps $\overline\Phi$ and $\Psi$ are inverse on generators, hence everywhere. 
\end{proof}

\begin{theorem}\label{thm:monoid-prime-explosion}
Assume that $R$ is a UFD with infinitely many pairwise nonassociate irreducibles. Let $\Pi$ be a fixed set of representatives of the irreducible classes. For $\Sigma\subseteq\Pi$, define
\[
\mathfrak p_\Sigma:=\{\tau,0_R\}\cup
\{x\in R\setminus\{0\}:p\mid x\text{ for some }p\in\Sigma\}.
\]
Then:
\begin{enumerate}[label=(\arabic*)]
    \item $\mathfrak p_\Sigma$ is a prime ideal of the commutative monoid $M_S$;
    \item for $\Sigma_1,\Sigma_2\subseteq\Pi$,
    \[
    \mathfrak p_{\Sigma_1}\subseteq \mathfrak p_{\Sigma_2}
    \iff
    \Sigma_1\subseteq\Sigma_2;
    \]
    \item $\Spec(M_S)$ has infinite Krull dimension.
\end{enumerate}
\end{theorem}

\begin{proof}
First, $\mathfrak p_\Sigma$ is an ideal: multiplying $\tau$ by anything gives $\tau$, multiplying $0_R$ by anything gives $0_R$ or $\tau$, and multiplying a supported element divisible by some $p\in\Sigma$ by another supported element preserves divisibility by $p$. It is proper because $1\notin\mathfrak p_\Sigma$.

Now suppose $ab\in\mathfrak p_\Sigma$. If $ab=\tau$, one factor is $\tau$. If $ab=0_R$, then $a,b\in R$, and since $R$ is a domain, one factor is $0_R$. If $ab\neq0_R$ and is divisible by some $p\in\Sigma$, then $a,b\in R\setminus\{0\}$ and irreducibles are prime in a UFD, so $p\mid a$ or $p\mid b$. Hence one factor lies in $\mathfrak p_\Sigma$. Thus $\mathfrak p_\Sigma$ is prime.

If $\Sigma_1\subseteq\Sigma_2$, then $\mathfrak p_{\Sigma_1}\subseteq\mathfrak p_{\Sigma_2}$ is immediate. Conversely, if $\mathfrak p_{\Sigma_1}\subseteq\mathfrak p_{\Sigma_2}$ and $p\in\Sigma_1$, then $p\in\mathfrak p_{\Sigma_2}$, so some $q\in\Sigma_2$ divides $p$. Since $p$ and $q$ are irreducible, they are associate; by the choice of $\Pi$, one has $q=p$. Thus $p\in\Sigma_2$.

Choose pairwise distinct $p_1,p_2,\dots\in\Pi$. Then
\[
\{\tau\}
\subsetneq
\mathfrak p_{\{p_1\}}
\subsetneq
\mathfrak p_{\{p_1,p_2\}}
\subsetneq\cdots
\]
is a chain of prime ideals of arbitrary finite length. 
\end{proof}

\begin{theorem}\label{thm:contracted-monoid-algebra}
Assume that $R$ is an integral domain. Let
\[
R^{\bullet}:=R\setminus\{0\}.
\]
Then
\[
\Z_0[M_S]:=\Z[M_S]/([\tau])\cong \Z\times \Z[R^{\bullet}].
\]
Consequently,
\[
\Spec(\Z_0[M_S])\cong \Spec(\Z)\sqcup \Spec(\Z[R^{\bullet}]).
\]
\end{theorem}

\begin{proof}
Killing $[\tau]$ removes the absorbing element, so
\[
\Z_0[M_S]\cong \Z[M],
\]
where $M=(R,\cdot)$.

Define
\[
\Phi:\Z[M]\to \Z\times \Z[R^{\bullet}]
\]
on basis elements by
\[
\Phi([0_R])=(1,0),
\qquad
\Phi([r])=(1,[r])
\quad(r\neq0).
\]
If one factor is $[0_R]$, multiplicativity is immediate. If $r,s\neq0$, then $rs\neq0$ because $R$ is a domain, and
\[
\Phi([r])\Phi([s])=(1,[r])(1,[s])=(1,[rs])=\Phi([rs]).
\]
Hence $\Phi$ is a ring homomorphism.

Define
\[
\Psi:\Z\times \Z[R^{\bullet}]\to \Z[M]
\]
by
\[
\Psi\!\left(c,\sum_{r\neq0}a_r[r]\right)=
\left(c-\sum_{r\neq0}a_r\right)[0_R]+\sum_{r\neq0}a_r[r].
\]
A direct computation shows $\Psi\Phi=\id$ and $\Phi\Psi=\id$. The spectrum statement follows from the elementary description of spectra of product rings. 
\end{proof}

\section{Semimodule classification: exact replacement of $R$-linear algebra}

\begin{lemma}\label{lem:epsilon-basic}
Let $M$ be a left $S$-semimodule and define
\[
\varepsilon_M(m):=em.
\]
Then:
\begin{enumerate}[label=(\arabic*)]
    \item $\varepsilon_M$ is additive;
    \item $\varepsilon_M^2=\varepsilon_M$;
    \item for every $m\in M$,
    \[
    m+\varepsilon_M(m)=m.
    \]
\end{enumerate}
\end{lemma}

\begin{proof}
Additivity and idempotence are immediate:
\[
\varepsilon_M(m+n)=e(m+n)=em+en,
\qquad
\varepsilon_M^2(m)=e(em)=e^2m=em.
\]
Finally,
\[
m+\varepsilon_M(m)=1\cdot m+e\cdot m=(1\oplus e)m=m,
\]
because $1\oplus e=1+0_R=1$ inside the supported copy of $R$. 
\end{proof}

\begin{proposition}\label{prop:support-semilattice}
Let
\[
L_M:=eM=\im(\varepsilon_M).
\]
Then $L_M$ is an idempotent commutative monoid under the ambient addition of $M$. In particular,
\[
l\vee l':=l+l'
\qquad(l,l'\in L_M)
\]
defines a join-semilattice structure on $L_M$ with least element $0_M$.
\end{proposition}

\begin{proof}
If $l=em$ and $l'=en$, then
\[
l+l'=e(m+n)\in L_M.
\]
Moreover,
\[
l+l=em+em=(e\oplus e)m=em=l,
\]
because $e\oplus e=e$. Thus $L_M$ is additively idempotent. Since $0_M=e0_M$, the least element is $0_M$. 
\end{proof}

\begin{theorem}\label{thm:fiber-modules}
For each $l\in L_M$, define the fiber
\[
M_l:=\{m\in M:\varepsilon_M(m)=l\}.
\]
Then:
\begin{enumerate}[label=(\arabic*)]
    \item $M_l$ is canonically an $R$-module;
    \item the zero element of the module $M_l$ is $l$ itself;
    \item if $l\le l'$ in $L_M$, then
    \[
    T_{l,l'}:M_l\to M_{l'},
    \qquad
    T_{l,l'}(m):=m+l'
    \]
    is a well-defined $R$-linear map;
    \item one has
    \[
    T_{l,l}=\id_{M_l},
    \qquad
    T_{l',l''}\circ T_{l,l'}=T_{l,l''}
    \quad(l\le l'\le l'').
    \]
\end{enumerate}
\end{theorem}

\begin{proof}
Fix $l\in L_M$.

If $m,n\in M_l$, then
\[
\varepsilon_M(m+n)=\varepsilon_M(m)+\varepsilon_M(n)=l+l=l,
\]
so $M_l$ is closed under addition.

Now $l$ is the additive identity on $M_l$: if $m\in M_l$, then by Lemma~\ref{lem:epsilon-basic},
\[
m+l=m+\varepsilon_M(m)=m.
\]
Also,
\[
m+(-1)m=(1\oplus(-1))m=0_Rm=em=l.
\]
Hence $(-1)m$ is the additive inverse of $m$ in the fiber. Thus $M_l$ is an abelian group.

If $r\in R$ and $m\in M_l$, then
\[
\varepsilon_M(rm)=e(rm)=(er)m=em=l,
\]
because $er=0_Rr=0_R=e$. Hence $rm\in M_l$. The module axioms are inherited from the semimodule axioms. Finally,
\[
0_Rm=em=l,
\]
which is the zero element of the module $M_l$. Thus $M_l$ is an $R$-module.

Assume $l\le l'$. If $m\in M_l$, then
\[
\varepsilon_M(m+l')=\varepsilon_M(m)+\varepsilon_M(l')=l+l'=l',
\]
so $m+l'\in M_{l'}$. Therefore $T_{l,l'}$ is well defined. It is additive because
\[
T_{l,l'}(m+n)=m+n+l'=(m+l')+(n+l'),
\]
using $l'+l'=l'$. It preserves the zero element because
\[
T_{l,l'}(l)=l+l'=l'.
\]
It is $R$-linear because if $r\in R$,
\[
T_{l,l'}(rm)=rm+l'=rm+rl'=r(m+l'),
\]
since every supported scalar acts trivially on $L_M$.

Finally,
\[
T_{l,l}(m)=m+l=m,
\]
and for $l\le l'\le l''$,
\[
T_{l',l''}(T_{l,l'}(m))=(m+l')+l''=m+(l'+l'')=m+l''=T_{l,l''}(m).
\]

\end{proof}

\begin{proposition}\label{prop:addition-via-join-fiber}
If $m\in M_l$ and $n\in M_{l'}$, then
\[
m+n
=
T_{l,l\vee l'}(m)
+_{M_{l\vee l'}}
T_{l',l\vee l'}(n),
\]
where $+_{M_{l\vee l'}}$ denotes the module addition inside the fiber $M_{l\vee l'}$.
\end{proposition}

\begin{proof}
The right-hand side equals
\[
(m+l\vee l')+(n+l\vee l').
\]
Since
\[
\varepsilon_M(m+n)=\varepsilon_M(m)+\varepsilon_M(n)=l\vee l',
\]
Lemma~\ref{lem:epsilon-basic} gives
\[
(m+n)+(l\vee l')=m+n.
\]
Using idempotence of $l\vee l'$, the expressions coincide. 
\end{proof}

\begin{definition}\label{def:support-diagram}
A \emph{support diagram over $R$} is a triple
\[
\mathcal A=\bigl(L,\{A_l\}_{l\in L},\{\rho_{l,l'}\}_{l\le l'}\bigr)
\]
where:
\begin{enumerate}[label=(\arabic*)]
    \item $L$ is a join-semilattice with least element $0$;
    \item each $A_l$ is an $R$-module;
    \item each $\rho_{l,l'}:A_l\to A_{l'}$ is $R$-linear;
    \item $\rho_{l,l}=\id_{A_l}$ and $\rho_{l',l''}\circ\rho_{l,l'}=\rho_{l,l''}$ whenever $l\le l'\le l''$.
\end{enumerate}
Morphisms are the obvious join-preserving natural transformations. The resulting category is denoted $\Diag_R$.
\end{definition}

\begin{theorem}\label{thm:reconstruct-semimodule}
Let
\[
\mathcal A=\bigl(L,\{A_l\}_{l\in L},\{\rho_{l,l'}\}_{l\le l'}\bigr)\in \Diag_R.
\]
Define
\[
\mathbb M(\mathcal A):=\bigsqcup_{l\in L}A_l.
\]
For $x\in A_l$ and $y\in A_{l'}$, define
\[
x\boxplus y:=\rho_{l,l\vee l'}(x)+\rho_{l',l\vee l'}(y)\in A_{l\vee l'}.
\]
Let $0$ denote the zero element of $A_0$. For $r\in R$, define
\[
r\cdot x:=rx\in A_l,
\qquad
\tau\cdot x:=0\in A_0.
\]
Then $\mathbb M(\mathcal A)$ is a left $S$-semimodule.
\end{theorem}

\begin{proof}
Commutativity of $\boxplus$ is immediate. To prove associativity, let
\[
x\in A_l,
\qquad
y\in A_m,
\qquad
z\in A_n,
\qquad
u:=l\vee m\vee n.
\]
Then
\[
(x\boxplus y)\boxplus z
=
\rho_{l\vee m,\nu}\bigl(\rho_{l,l\vee m}(x)+\rho_{m,l\vee m}(y)\bigr)+\rho_{n,\nu}(z).
\]
By linearity and functoriality of the $\rho$'s this equals
\[
\rho_{l,\nu}(x)+\rho_{m,\nu}(y)+\rho_{n,\nu}(z),
\]
which is symmetric in the three terms. Hence $\boxplus$ is associative.

If $x\in A_l$, then
\[
x\boxplus 0=\rho_{l,l}(x)+\rho_{0,l}(0)=x+0_{A_l}=x,
\]
and similarly $0\boxplus x=x$.

Now let $r,s\in R$. If $x\in A_l$, then
\[
(r+s)\cdot x=(r+s)x=rx+sx=(r\cdot x)\boxplus (s\cdot x),
\]
because both terms lie in the same fiber $A_l$. Also
\[
(rs)\cdot x=r\cdot(s\cdot x),
\qquad
1\cdot x=x.
\]
If $x\in A_l$ and $y\in A_{l'}$, then
\[
r\cdot(x\boxplus y)
=
r\bigl(\rho_{l,l\vee l'}(x)+\rho_{l',l\vee l'}(y)\bigr)
=
\rho_{l,l\vee l'}(rx)+\rho_{l',l\vee l'}(ry)
=
(r\cdot x)\boxplus(r\cdot y).
\]
Finally,
\[
\tau\cdot(x\boxplus y)=0=(\tau\cdot x)\boxplus(\tau\cdot y),
\]
and
\[
(r\oplus\tau)\cdot x=r\cdot x=(r\cdot x)\boxplus(\tau\cdot x).
\]
Hence all semimodule axioms hold. 
\end{proof}

\begin{theorem}\label{thm:equivalence-semimodules-diagrams}
The assignments
\[
M\longmapsto\bigl(L_M,\{M_l\}_{l\in L_M},\{T_{l,l'}\}_{l\le l'}\bigr)
\]
and
\[
\mathcal A\longmapsto \mathbb M(\mathcal A)
\]
extend to mutually quasi-inverse equivalences
\[
\boxed{
S\text{-}\SMod\simeq \Diag_R.
}
\]
\end{theorem}

\begin{proof}
The extraction of a support diagram from an $S$-semimodule is justified by Theorem~\ref{thm:fiber-modules} and Proposition~\ref{prop:addition-via-join-fiber}. Conversely, Theorem~\ref{thm:reconstruct-semimodule} constructs an $S$-semimodule from any support diagram.

These constructions are functorial. If $f:M\to N$ is $S$-linear, then
\[
\varepsilon_N(f(m))=ef(m)=f(em)=f(\varepsilon_M(m)),
\]
so $f$ carries each fiber $M_l$ into the corresponding fiber of $N$, and the resulting maps commute with transport.

If $(\phi,\eta):\mathcal A\to\mathcal A'$ is a morphism of support diagrams, then the map on the disjoint unions is additive because naturality gives
\[
\eta_{l\vee l'}(\rho_{l,l\vee l'}(x)+\rho_{l',l\vee l'}(y))
=
\rho'_{\phi(l),\phi(l)\vee\phi(l')}(\eta_l(x))
+
\rho'_{\phi(l'),\phi(l)\vee\phi(l')}(\eta_{l'}(y)).
\]

Starting from an $S$-semimodule $M$, the reconstructed semimodule has the same underlying set, because the fibers partition $M$. Proposition~\ref{prop:addition-via-join-fiber} shows that the reconstructed addition agrees with the original one. Scalar compatibility is immediate.

Starting from a support diagram $\mathcal A$, the reconstructed semimodule has support semilattice identified with $L$ via the zero elements of the fibers, and the extracted fibers and transport maps are the original ones by construction. 
\end{proof}

\begin{corollary}\label{cor:cancellative-category}
Restriction along $p_R$ induces an equivalence
\[
R\text{-}\Mod\simeq S\text{-}\SMod^{\mathrm{canc}},
\]
where the right-hand side is the full subcategory of cancellative $S$-semimodules.
\end{corollary}

\begin{proof}
If $M$ is cancellative, Lemma~\ref{lem:epsilon-basic} gives
\[
m+\varepsilon_M(m)=m.
\]
Cancellation implies $\varepsilon_M(m)=0_M$ for all $m$. Hence $L_M$ is the one-point semilattice, so Theorem~\ref{thm:equivalence-semimodules-diagrams} identifies $M$ with a single $R$-module.

Conversely, any $R$-module $A$ becomes an $S$-semimodule by
\[
r\cdot a:=ra,
\qquad
\tau\cdot a:=0,
\]
and this semimodule is cancellative. 
\end{proof}

\begin{corollary}\label{cor:adjunctions}
The functor
\[
M\longmapsto eM
\]
is left adjoint to restriction along $\chi_R:S\to\B$, and the functor
\[
M\longmapsto K_M:=\ker(\varepsilon_M)
\]
is right adjoint to restriction along $p_R:S\to R$.
\end{corollary}

\begin{proof}
Let $N$ be a $\B$-semimodule, regarded as an $S$-semimodule via $\chi_R$. If $f:M\to N$ is $S$-linear, then for every $m\in M$,
\[
f(m)=e\cdot f(m)=f(em),
\]
because $e$ is supported and so acts as $1$ on $N$. Thus $f$ factors uniquely through $\varepsilon_M:M\to eM$, giving
\[
\Hom_S(M,\chi_R^*N)\cong \Hom_{\B}(eM,N).
\]

Now let $A$ be an $R$-module, viewed as an $S$-semimodule via $p_R$. If $g:A\to M$ is $S$-linear, then
\[
eg(a)=g(ea)=g(0)=0_M,
\]
so $g(A)\subseteq K_M$. Conversely, any $R$-linear map $A\to K_M$ is automatically $S$-linear. Hence
\[
\Hom_S(p_R^*A,M)\cong \Hom_R(A,K_M).
\]

\end{proof}

\begin{theorem}\label{thm:free-semimodule-boolean-cube}
Let $F_n:=S^n$. Then the support diagram of $F_n$ is the Boolean lattice
\[
\cP(\{1,\dots,n\})
\]
with fiber over $I\subseteq\{1,\dots,n\}$ canonically equal to $R^I$, and transport along inclusions $I\subseteq J$ given by extension by zero in the coordinates of $J\setminus I$.
\end{theorem}

\begin{proof}
The support of a vector $x=(x_1,\dots,x_n)\in S^n$ is determined exactly by the set of indices $j$ for which $x_j\in R$. Thus
\[
L_{F_n}\cong \cP(\{1,\dots,n\}),
\]
with join equal to union.

Fix $I\subseteq\{1,\dots,n\}$. The fiber over $I$ consists of vectors whose coordinates in $I$ are supported and whose coordinates outside $I$ are $\tau$. Projection to the supported coordinates identifies that fiber with $R^I$. If $I\subseteq J$, adding the support idempotent of $J$ inserts zeros in the new coordinates $J\setminus I$. 
\end{proof}

\begin{theorem}\label{thm:subsemimodule-classification}
Let $M$ be an $S$-semimodule. A subsemimodule $N\subseteq M$ is equivalent to the data of:
\begin{enumerate}[label=(\arabic*)]
    \item a join-subsemilattice with zero $L'\subseteq L_M$;
    \item for each $l\in L'$, an $R$-submodule $N_l\subseteq M_l$ containing the zero element $l$ of the fiber;
    \item compatibility with transport:
    \[
    T_{l,l'}(N_l)\subseteq N_{l'}
    \qquad(l\le l'\text{ in }L').
    \]
\end{enumerate}
Under this correspondence,
\[
N=\bigsqcup_{l\in L'}N_l.
\]
\end{theorem}

\begin{proof}
If $N\subseteq M$ is a subsemimodule, then $eN\subseteq N$, so
\[
L_N:=eN\subseteq L_M.
\]
If $l,l'\in L_N$, choose $x,y\in N$ with $ex=l$ and $ey=l'$. Then
\[
l\vee l'=e(x+y)\in eN=L_N,
\]
so $L_N$ is a join-subsemilattice with zero.

For $l\in L_N$, let
\[
N_l:=N\cap M_l.
\]
Because $N$ is stable under addition and under the action of $R\subset S$, each $N_l$ is an $R$-submodule of $M_l$. If $l\le l'$ and $x\in N_l$, then
\[
T_{l,l'}(x)=x+l'\in N,
\]
because both $x$ and $l'$ lie in $N$. Hence $T_{l,l'}(N_l)\subseteq N_{l'}$.

Conversely, given such data, define
\[
N:=\bigsqcup_{l\in L'}N_l.
\]
Then $0_M\in N$. If $x\in N_l$ and $y\in N_{l'}$, then by Proposition~\ref{prop:addition-via-join-fiber},
\[
x+y
=
T_{l,l\vee l'}(x)+_{M_{l\vee l'}}T_{l',l\vee l'}(y)
\in N_{l\vee l'}.
\]
Thus $N$ is closed under addition. It is stable under $R$ because each $N_l$ is an $R$-submodule, and under $\tau$ because $\tau x=0_M$. Hence $N$ is an $S$-subsemimodule. 
\end{proof}

\begin{theorem}\label{thm:finite-generation}
Assume that $M$ is generated as an $S$-semimodule by $m_1,\dots,m_n$. Set
\[
l_i:=\varepsilon_M(m_i)\in L_M.
\]
Then:
\begin{enumerate}[label=(\arabic*)]
    \item $L_M$ is the join-subsemilattice generated by $l_1,\dots,l_n$; in particular,
    \[
    |L_M|\le 2^n;
    \]
    \item for each $l\in L_M$, the fiber $M_l$ is generated as an $R$-module by the transported elements
    \[
    T_{l_i,l}(m_i)
    \qquad
    (1\le i\le n,\ l_i\le l).
    \]
\end{enumerate}
\end{theorem}

\begin{proof}
Every element of $M$ is a finite sum of terms $r_jm_{i_j}$ with $r_j\in R$. Since supported scalars preserve support,
\[
\varepsilon_M(r_jm_{i_j})=\varepsilon_M(m_{i_j})=l_{i_j}.
\]
Therefore
\[
\varepsilon_M\!\left(\sum_j r_jm_{i_j}\right)=\bigvee_j l_{i_j}.
\]
Hence every support element is a join of some subset of $\{l_1,\dots,l_n\}$, and every such join occurs as the support of the corresponding sum of generators. This proves (1).

Now let $m\in M_l$ and write
\[
m=\sum_j r_jm_{i_j}.
\]
Then
\[
l=\varepsilon_M(m)=\bigvee_j l_{i_j},
\]
so each $l_{i_j}\le l$. Therefore each $T_{l_{i_j},l}(m_{i_j})$ is defined. Since transport is $R$-linear,
\[
\sum_j r_jT_{l_{i_j},l}(m_{i_j})
=
\sum_j T_{l_{i_j},l}(r_jm_{i_j})
=
\sum_j (r_jm_{i_j}+l)
=
m+l
=
m,
\]
because $l$ is the zero of the module $M_l$. 
\end{proof}

\section{Conservative ideal, class, and factorization phenomena}

\begin{theorem}\label{thm:ideal-multiplication-classical}
Assume the baseline classification of ideals,
\[
I_J:=J\cup\{\tau\}
\qquad(J\triangleleft R).
\]
Then for all ideals $J,K\triangleleft R$,
\[
I_JI_K=I_{JK}.
\]
If $R$ is an integral domain, then the ideal class monoid of the nontrivial ideals of $G(R)$ is canonically isomorphic to the ordinary ideal class monoid of $R$.
\end{theorem}

\begin{proof}
Every product of an element of $I_J$ with an element of $I_K$ is either $\tau$ or an element of $JK$, so $I_JI_K\subseteq I_{JK}$. Conversely, $\tau\in I_JI_K$, and every element of $JK$ is a finite sum of products $ab$ with $a\in J$ and $b\in K$, hence belongs to $I_JI_K$. Thus $I_JI_K=I_{JK}$.

Now assume $R$ is a domain. For nonzero ideals $J_1,J_2\subset R$, classical ideal-class equivalence is
\[
J_1\sim_R J_2
\iff
\exists a,b\in R\setminus\{0\}:aJ_1=bJ_2.
\]
For the corresponding $G(R)$-ideals,
\[
I_{J_1}\sim_G I_{J_2}
\iff
\exists a,b\in R\setminus\{0\}:aI_{J_1}=bI_{J_2}.
\]
But
\[
aI_J=aJ\cup\{\tau\}.
\]
Therefore
\[
aI_{J_1}=bI_{J_2}
\iff
aJ_1=bJ_2.
\]
Hence the class monoids coincide. 
\end{proof}

\begin{proposition}\label{prop:separated-weight-nonmultiplicative}
Assume that every nonzero quotient $R/J$ is finite, and write
\[
N(J):=|R/J|.
\]
Then the separated-sheet weight
\[
w(J):=N(J)+1
\]
is not multiplicative on nonzero ideals. More precisely,
\[
w(JK)=N(J)N(K)+1,
\qquad
w(J)w(K)=N(J)N(K)+N(J)+N(K)+1.
\]
Consequently the separated congruence zeta sheet
\[
\zeta^{\mathrm{sep}}_{G(R)}(s):=\sum_{0\neq J}(N(J)+1)^{-s}
\]
is not an Euler product indexed by prime ideals.
\end{proposition}

\begin{proof}
The displayed identities are immediate from the definition. Their difference is $N(J)+N(K)$, which is strictly positive. Therefore $w$ is not multiplicative. 
\end{proof}

\begin{proposition}\label{prop:same-abscissa}
Under the same finiteness hypothesis, the classical ideal zeta sheet
\[
\zeta^{\mathrm{idl}}_{G(R)}(s):=\sum_{0\neq J}N(J)^{-s}
\]
and the separated sheet $\zeta^{\mathrm{sep}}_{G(R)}(s)$ have the same abscissa of absolute convergence.
\end{proposition}

\begin{proof}
For every nonzero ideal $J$ and every $\sigma>0$,
\[
N(J)\le N(J)+1\le 2N(J),
\]
whence
\[
2^{-\sigma}N(J)^{-\sigma}\le (N(J)+1)^{-\sigma}\le N(J)^{-\sigma}.
\]
The conclusion follows by comparison. 
\end{proof}

\begin{proposition}\label{prop:units-irreducibles}
Assume that $R$ is an integral domain. Then:
\begin{enumerate}[label=(\arabic*)]
    \item the units of $G(R)$ are exactly the units of $R$;
    \item every nonzero supported irreducible of $G(R)$ is exactly an irreducible of $R$, and conversely.
\end{enumerate}
\end{proposition}

\begin{proof}
If $u\in G(R)$ is a unit, then $u\neq \tau$ because $\tau x=\tau\neq1$. Hence $u\in R$, and its inverse is also supported, so $u$ is a unit of $R$. The converse is immediate.

Let $x\in R\setminus\{0\}$. If $x=ab$ in $G(R)$, then neither factor can be $\tau$, since otherwise $ab=\tau$. Hence $a,b\in R$, so factorization of $x$ in $G(R)$ is exactly factorization in $R$. 
\end{proof}

\begin{theorem}\label{thm:supported-zero-prime-reducible}
Assume that $R$ is an integral domain. The supported zero $e=0_R\in G(R)$ is reducible and prime in the divisibility sense.
\end{theorem}

\begin{proof}
It is reducible because
\[
e=e\cdot e,
\]
and $e$ is not a unit.

Now suppose $e\mid xy$. Then
\[
xy=e\cdot c
\]
for some $c\in G(R)$. If $c=\tau$, then $xy=\tau$. If $c\in R$, then $xy=e$. Hence
\[
xy\in\{e,\tau\}.
\]
If both $x$ and $y$ are supported nonzero elements, then $xy$ is again a supported nonzero element of $R$, impossible. Hence one of $x,y$ lies in $\{e,\tau\}$, so $e$ divides one of them. 
\end{proof}

\begin{proposition}\label{prop:scalar-norm-collapse}
Assume that $R$ is an integral domain with a multiplicative ideal norm $N$ on nonzero ideals, and assume that there exists a proper nonzero ideal $J$ with $N(J)>1$. Let
\[
I_{(0)}:=\{e,\tau\}.
\]
If
\[
h:\{\text{ideals of }G(R)\}\to [0,\infty)
\]
satisfies
\[
h(I_1I_2)=h(I_1)h(I_2)
\]
for all ideals and
\[
h(I_J)=N(J)
\]
for every nonzero ideal $J\triangleleft R$, then
\[
h(I_{(0)})=0,
\qquad
h(\{\tau\})=0.
\]
\end{proposition}

\begin{proof}
Since $I_{(0)}I_J=I_{(0)}$, one has
\[
h(I_{(0)})=h(I_{(0)})h(I_J)=h(I_{(0)})N(J).
\]
Because $N(J)\neq1$, it follows that $h(I_{(0)})=0$. Similarly,
\[
\{\tau\}I_J=\{\tau\}
\]
implies
\[
h(\{\tau\})=h(\{\tau\})N(J),
\]
so $h(\{\tau\})=0$. 
\end{proof}

\section{Universal all-orders deformation of the shifted zeta sheet}

\begin{definition}\label{def:Lt}
Let
\[
L(s)=\sum_{n\ge1}a_n n^{-s}
\]
be a Dirichlet series with abscissa of absolute convergence $\sigma_a$. For $|t|<1$, define
\[
L_t(s):=\sum_{n\ge1}a_n(n+t)^{-s}.
\]
Let
\[
(Tf)(s):=f(s+1),
\qquad
D:=sT.
\]
\end{definition}

\begin{theorem}\label{thm:universal-taylor-deformation}
For $|t|<1$ and $\Re(s)>\sigma_a$,
\[
\boxed{
L_t(s)=\sum_{m\ge0}(-1)^m\frac{(s)_m}{m!}t^mL(s+m).
}
\]
Equivalently,
\[
\boxed{
L_t=e^{-tD}L.
}
\]
\end{theorem}

\begin{proof}
For $|u|<1$,
\[
(1+u)^{-s}=\sum_{m\ge0}(-1)^m\frac{(s)_m}{m!}u^m.
\]
Apply this with $u=t/n$. Since $|t|<1$, one has $|t/n|<1$ for every $n\ge1$, so
\[
(n+t)^{-s}
=
n^{-s}\Bigl(1+\frac{t}{n}\Bigr)^{-s}
=
\sum_{m\ge0}(-1)^m\frac{(s)_m}{m!}t^m n^{-s-m}.
\]
Multiplying by $a_n$ and summing over $n$ gives
\[
L_t(s)
=
\sum_{m\ge0}(-1)^m\frac{(s)_m}{m!}t^m\sum_{n\ge1}a_n n^{-s-m}
=
\sum_{m\ge0}(-1)^m\frac{(s)_m}{m!}t^mL(s+m).
\]
Also,
\[
D^mL(s)=(s)_mL(s+m)
\]
by induction on $m$. Hence
\[
e^{-tD}L=
\sum_{m\ge0}\frac{(-t)^m}{m!}D^mL
=
\sum_{m\ge0}(-1)^m\frac{(s)_m}{m!}t^mL(s+m)=L_t.
\]

\end{proof}

\begin{corollary}\label{cor:jets}
For every integer $n\ge0$,
\[
\boxed{
\partial_t^{\,n}L_t(s)=(-1)^n(s)_nL_t(s+n).
}
\]
In particular,
\[
\boxed{
\partial_t^{\,n}L_t(s)\big|_{t=0}=(-1)^n(s)_nL(s+n).
}
\]
\end{corollary}

\begin{proof}
Differentiate termwise:
\[
\partial_t^{\,n}(m+t)^{-s}=(-1)^n(s)_n(m+t)^{-s-n}.
\]
Summing over $m$ yields the formula. 
\end{proof}

\begin{corollary}\label{cor:semigroup-law}
If $|t_0|<1$ and $|h|<1-|t_0|$, then
\[
L_{t_0+h}=e^{-hD}L_{t_0}.
\]
\end{corollary}

\begin{proof}
Apply Theorem~\ref{thm:universal-taylor-deformation} to the Dirichlet series
\[
s\longmapsto L_{t_0}(s)=\sum_{n\ge1}a_n(n+t_0)^{-s}.
\]

\end{proof}

\begin{theorem}\label{thm:meromorphic-continuation-shifted-tower}
Assume that $L$ extends meromorphically to $\C$. Let $N\ge1$ and $|t|<1$. Then on the half-plane $\Re(s)>\sigma_a-N$,
\[
\boxed{
L_t(s)=
\sum_{m=0}^{N-1}(-1)^m\frac{(s)_m}{m!}t^mL(s+m)
+
H_{N,t}(s),
}
\]
where $H_{N,t}(s)$ is holomorphic on $\Re(s)>\sigma_a-N$.
\end{theorem}

\begin{proof}
Apply Taylor's formula with integral remainder to
\[
f_s(u):=(1+u)^{-s}
\]
at $u=0$, then substitute $u=t/n$. One obtains
\[
(n+t)^{-s}
=
\sum_{m=0}^{N-1}(-1)^m\frac{(s)_m}{m!}t^m n^{-s-m}
+
t^N n^{-s-N}R_N(t/n,s),
\]
where $R_N(t/n,s)$ is holomorphic in $s$ and uniformly bounded on compact subsets of $\Re(s)>\sigma_a-N$. Multiplying by $a_n$ and summing over $n$ gives the stated decomposition. The remainder series converges absolutely and locally uniformly on $\Re(s)>\sigma_a-N$, hence defines a holomorphic function there. 
\end{proof}

\begin{corollary}\label{cor:pole-rigidity}
If $L$ has a unique simple pole at $s=1$, then every $L_t$ has the same unique simple pole at $s=1$, with the same residue:
\[
\operatorname*{Res}_{s=1}L_t(s)=\operatorname*{Res}_{s=1}L(s).
\]
\end{corollary}

\begin{proof}
In Theorem~\ref{thm:meromorphic-continuation-shifted-tower}, the only term that can contribute a pole at $s=1$ is $L(s)$ itself, since $L(s+m)$ is regular there for all $m\ge1$. 
\end{proof}

\begin{definition}\label{def:cumulants}
Assume $L_t(s)\neq0$. Define
\[
A_n(t;s):=(s)_n\frac{L_t(s+n)}{L_t(s)}
\qquad(n\ge1),
\]
and define the logarithmic cumulants
\[
\kappa_n(t;s):=\partial_t^{\,n}\log L_t(s).
\]
\end{definition}

\begin{theorem}\label{thm:cumulant-partition-formula}
For every $n\ge1$,
\[
\boxed{
\kappa_n(t;s)
=
(-1)^n
\sum_{\pi\in\Pi_n}
(-1)^{|\pi|-1}(|\pi|-1)!
\prod_{B\in\pi}A_{|B|}(t;s),
}
\]
where $\Pi_n$ is the set of set partitions of $\{1,\dots,n\}$.
\end{theorem}

\begin{proof}
Set
\[
m_n(t;s):=\frac{\partial_t^{\,n}L_t(s)}{L_t(s)}.
\]
By Corollary~\ref{cor:jets},
\[
m_n(t;s)=(-1)^nA_n(t;s).
\]
Now form the exponential generating series
\[
M(z):=\sum_{n\ge0}m_n\frac{z^n}{n!},
\qquad
K(z):=\sum_{n\ge1}\kappa_n\frac{z^n}{n!}.
\]
Taylor expansion yields
\[
L_{t+z}(s)=L_t(s)M(z),
\]
whereas the definition of $\kappa_n$ gives
\[
\log L_{t+z}(s)-\log L_t(s)=K(z),
\]
so
\[
L_{t+z}(s)=L_t(s)e^{K(z)}.
\]
Hence
\[
M(z)=e^{K(z)},
\qquad
K(z)=\log M(z).
\]
The standard moment--cumulant formula for exponential generating series implies
\[
\kappa_n=
\sum_{\pi\in\Pi_n}
(-1)^{|\pi|-1}(|\pi|-1)!
\prod_{B\in\pi}m_{|B|}.
\]
Substituting $m_{|B|}=(-1)^{|B|}A_{|B|}$ and using $\sum_{B\in\pi}|B|=n$ yields the displayed formula. 
\end{proof}

\begin{corollary}\label{cor:first-four-cumulants}
The first four cumulants are
\[
\kappa_1=-A_1,
\]
\[
\kappa_2=A_2-A_1^2,
\]
\[
\kappa_3=-A_3+3A_1A_2-2A_1^3,
\]
\[
\kappa_4=A_4-4A_1A_3-3A_2^2+12A_1^2A_2-6A_1^4.
\]
\end{corollary}

\begin{proof}
Expand the partition formula of Theorem~\ref{thm:cumulant-partition-formula} for $n=1,2,3,4$. 
\end{proof}

\begin{theorem}\label{thm:arithmetic-kernels-riemann}
For every integer $r\ge1$, the ratio
\[
Q_r(s):=\frac{\zeta(s+r)}{\zeta(s)}
\]
has a Dirichlet expansion
\[
Q_r(s)=\sum_{n\ge1}\alpha_r(n)n^{-s},
\]
where $\alpha_r$ is multiplicative and
\[
\alpha_r(p^k)=-\frac{p^r-1}{p^{rk}}
\qquad(k\ge1).
\]
Equivalently,
\[
\boxed{
\alpha_r(n)=(-1)^{\omega(n)}\frac{\prod_{p\mid n}(p^r-1)}{n^r}.
}
\]
\end{theorem}

\begin{proof}
Using the Euler product for $\zeta$,
\[
Q_r(s)=\prod_p\frac{1-p^{-s}}{1-p^{-s-r}}.
\]
At a fixed prime $p$, set $u=p^{-s}$. Then
\[
\frac{1-u}{1-u/p^r}
=
(1-u)\sum_{k\ge0}p^{-rk}u^k
=
1+\sum_{k\ge1}(p^{-rk}-p^{-r(k-1)})u^k
=
1-\sum_{k\ge1}\frac{p^r-1}{p^{rk}}u^k.
\]
Hence
\[
\alpha_r(p^0)=1,
\qquad
\alpha_r(p^k)=-\frac{p^r-1}{p^{rk}}
\quad(k\ge1).
\]
Multiplicativity follows from the Euler product, and the explicit formula for general $n$ is the product over the primes dividing $n$. 
\end{proof}

\begin{corollary}\label{cor:first-log-jet-explicit}
One has
\[
\partial_t\log\zeta_t(s)\Big|_{t=0}
=
-s\sum_{n\ge1}\alpha_1(n)n^{-s},
\]
where
\[
\alpha_1(n)=(-1)^{\omega(n)}\frac{\varphi(\rad(n))}{n}.
\]
Moreover,
\[
\partial_t^2\log\zeta_t(s)\Big|_{t=0}
=
s(s+1)\sum_{n\ge1}\alpha_2(n)n^{-s}
-
s^2\left(\sum_{n\ge1}\alpha_1(n)n^{-s}\right)^2.
\]
Hence the second logarithmic jet is already not primewise local.
\end{corollary}

\begin{proof}
By Theorem~\ref{thm:arithmetic-kernels-riemann},
\[
\frac{\zeta(s+1)}{\zeta(s)}=\sum_{n\ge1}\alpha_1(n)n^{-s},
\qquad
\frac{\zeta(s+2)}{\zeta(s)}=\sum_{n\ge1}\alpha_2(n)n^{-s}.
\]
Now Corollary~\ref{cor:first-four-cumulants} evaluated at $t=0$ gives
\[
\partial_t\log\zeta_t(s)\Big|_{t=0}=-A_1(0;s)=-s\frac{\zeta(s+1)}{\zeta(s)},
\]
and
\[
\partial_t^2\log\zeta_t(s)\Big|_{t=0}=A_2(0;s)-A_1(0;s)^2.
\]
Substituting the explicit Dirichlet expansions yields the formulas. The identity
\[
\prod_{p\mid n}(p-1)=\varphi(\rad(n))
\]
gives the expression for $\alpha_1(n)$. The square term involves Dirichlet convolution and therefore mixes prime supports. 
\end{proof}

\begin{theorem}\label{thm:non-eulerian-nonintegral-t}
Let
\[
L(s)=\sum_{n\ge1}a_n n^{-s}
\]
with $a_1\neq0$. If $0<t<1$, then $L_t$ is not an ordinary Dirichlet series in the variable $n^{-s}$. In particular, $L_t$ admits no classical Euler product.
\end{theorem}

\begin{proof}
Assume that
\[
L_t(s)=\sum_{n\ge1}b_n n^{-s}
\]
for $\Re(s)$ sufficiently large. Let $n_0$ be the least integer with $b_{n_0}\neq0$. Then as $\sigma\to+\infty$,
\[
\sum_{n\ge1}b_n n^{-\sigma}=b_{n_0}n_0^{-\sigma}(1+o(1)).
\]
Hence the leading exponential scale is $n_0^{-\sigma}$.

On the other hand,
\[
L_t(\sigma)=\sum_{n\ge1}a_n(n+t)^{-\sigma}=a_1(1+t)^{-\sigma}(1+o(1)),
\qquad \sigma\to+\infty,
\]
because the first term dominates all later terms. Hence the leading exponential scale is $(1+t)^{-\sigma}$, so one must have $n_0=1+t$. But $1+t\in(1,2)$ is not an integer. Contradiction. 
\end{proof}

\section{Hurwitz specialization and exact arithmetic consequences}

\begin{definition}\label{def:hurwitz-specialization}
Set
\[
F_t(s):=\sum_{n\ge1}(n+t)^{-s}=\zeta(s,1+t),
\qquad
D_t(s):=\zeta(s)-F_t(s).
\]
\end{definition}

\begin{proposition}\label{prop:hurwitz-standard-package}
The following identities hold.
\begin{enumerate}[label=(\arabic*)]
    \item Near $s=1$,
    \[
    F_t(s)=\frac1{s-1}-\psi(1+t)+O(s-1).
    \]
    \item For every integer $m\ge0$,
    \[
    F_t(-m)=-\frac{B_{m+1}(1+t)}{m+1}.
    \]
    \item One has
    \[
    F_t'(0)=\log\Gamma(1+t)-\frac12\log(2\pi).
    \]
\end{enumerate}
\end{proposition}

\begin{proof}
These are exactly the standard identities for the Hurwitz zeta function $\zeta(s,a)$ at $a=1+t$. See the references at the end. 
\end{proof}

\begin{corollary}\label{cor:defect-entire-bernoulli-gamma}
The defect $D_t(s)=\zeta(s)-F_t(s)$ is entire and satisfies
\[
D_t(-m)=\frac{B_{m+1}(1+t)-B_{m+1}}{m+1}
\qquad(m\ge0),
\]
and
\[
D_t'(0)=-\log\Gamma(1+t),
\qquad
e^{-D_t'(0)}=\Gamma(1+t).
\]
\end{corollary}

\begin{proof}
Both $\zeta(s)$ and $F_t(s)$ have a unique simple pole at $s=1$ of residue $1$, so their difference is entire. The formulas follow by subtracting the classical identities
\[
\zeta(-m)=-\frac{B_{m+1}}{m+1},
\qquad
\zeta'(0)=-\frac12\log(2\pi),
\]
from Proposition~\ref{prop:hurwitz-standard-package}. 
\end{proof}

\begin{proposition}\label{prop:integer-shift-coefficient-surgery}
For every integer $q\ge0$,
\[
F_q(s)=\zeta(s)-\sum_{n=1}^{q}n^{-s},
\qquad
D_q(s)=\sum_{n=1}^{q}n^{-s}.
\]
\end{proposition}

\begin{proof}
By the defining series of the Hurwitz zeta,
\[
F_q(s)=\zeta(s,1+q)=\sum_{m\ge0}(m+1+q)^{-s}=\sum_{n\ge q+1}n^{-s}.
\]
Subtract this from $\zeta(s)=\sum_{n\ge1}n^{-s}$. 
\end{proof}

\begin{corollary}\label{cor:faulhaber-factorial}
For every integer $q\ge0$ and every integer $m\ge0$,
\[
D_q(-m)=\sum_{n=1}^{q}n^m.
\]
Moreover,
\[
e^{-D_q'(0)}=q!.
\]
\end{corollary}

\begin{proof}
Evaluate Proposition~\ref{prop:integer-shift-coefficient-surgery} at $s=-m$. The factorial identity is Corollary~\ref{cor:defect-entire-bernoulli-gamma} with $t=q$, since $\Gamma(1+q)=q!$. 
\end{proof}

\section{The endpoint $t=1$: exact Stone defect and coprime kernel}

To avoid conflict with the supported zero $e=0_R\in S$, denote the characteristic functions on $\N$ by
\[
\mathbf u:=\one_{\{1\}},
\qquad
\mathbf v:=\one_{\N_{\ge2}}.
\]
Let
\[
X:=\{\xi_1,\xi_\ast\}
\]
be the Stone space of the clopen partition
\[
\N=\{1\}\sqcup\N_{\ge2}.
\]
The multiplication on $\N$ induces the semigroup law
\[
\xi_1\xi_1=\xi_1,
\qquad
\xi_1\xi_\ast=\xi_\ast\xi_1=\xi_\ast,
\qquad
\xi_\ast\xi_\ast=\xi_\ast.
\]
Let $\Delta$ be the pullback comultiplication on functions on $X$:
\[
(\Delta g)(x,y):=g(xy).
\]
Let
\[
\mathcal D(a)(s):=\sum_{n\ge1}a(n)n^{-s}
\]
be the Dirichlet transform.

\begin{theorem}\label{thm:F1-defect-and-boundary}
For
\[
F_1(s)=\zeta(s)-1=\sum_{n\ge1}b(n)n^{-s},
\qquad
b(1)=0,
\quad
b(n)=1\ (n\ge2),
\]
the coprime multiplicativity defect
\[
\Delta_{F_1}(m,n):=b(mn)-b(m)b(n)
\qquad((m,n)=1)
\]
is given by
\[
\Delta_{F_1}(m,n)=
\begin{cases}
1,&\text{if exactly one of }m,n\text{ equals }1,\\[1mm]
0,&\text{otherwise.}
\end{cases}
\]
Equivalently,
\[
\sum_{(m,n)=1}\Delta_{F_1}(m,n)m^{-s}n^{-t}=\zeta(s)+\zeta(t)-2.
\]
Moreover,
\[
\mathcal D(\mathbf u)=1,
\qquad
\mathcal D(\mathbf v)=\zeta(s)-1,
\]
and
\[
\Delta(\mathbf u)=\mathbf u\otimes\mathbf u,
\qquad
\Delta(\mathbf v)=\mathbf v\otimes\mathbf u+\mathbf u\otimes\mathbf v+\mathbf v\otimes\mathbf v.
\]
Thus
\[
\Delta(\mathbf v)-\mathbf v\otimes\mathbf v
=
\mathbf v\otimes\mathbf u+\mathbf u\otimes\mathbf v,
\]
so the non-Eulerianity of $F_1$ is supported exactly on the unit boundary.
\end{theorem}

\begin{proof}
The coefficient formula for $F_1$ is immediate. If both $m,n>1$, then $b(m)=b(n)=b(mn)=1$, so the defect is $0$. If $m=n=1$, then the defect is $0$. If exactly one of $m,n$ equals $1$, the defect is $1$. Hence the displayed formula holds, and summing over coprime pairs gives
\[
\sum_{n\ge2}n^{-t}+\sum_{m\ge2}m^{-s}=\zeta(s)+\zeta(t)-2.
\]
The formulas for $\mathcal D(\mathbf u)$ and $\mathcal D(\mathbf v)$ are immediate. The comultiplication formulas are obtained by evaluation on the four pairs $(\xi_1,\xi_1)$, $(\xi_1,\xi_\ast)$, $(\xi_\ast,\xi_1)$, $(\xi_\ast,\xi_\ast)$. Subtracting gives the final identity. 
\end{proof}

\begin{theorem}\label{thm:two-sheet-defect}
For
\[
\mathcal Z_{1,u}(s):=(1+u)\zeta(s)-u,
\qquad
\lambda:=1+u,
\]
the coefficient function is
\[
c_u(1)=1,
\qquad
c_u(n)=\lambda
\quad(n\ge2).
\]
Hence the coprime multiplicativity defect
\[
\Delta_u(m,n):=c_u(mn)-c_u(m)c_u(n)
\qquad((m,n)=1)
\]
is given by
\[
\Delta_u(m,n)=
\begin{cases}
0,&\min(m,n)=1,\\[1mm]
\lambda-\lambda^2=-u(1+u),&m,n>1.
\end{cases}
\]
Equivalently, if
\[
a_\lambda:=\mathbf u+\lambda\mathbf v,
\]
then
\[
\Delta(a_\lambda)-a_\lambda\otimes a_\lambda
=
(\lambda-\lambda^2)\,\mathbf v\otimes\mathbf v.
\]
Thus the scalar
\[
\delta(\lambda):=\lambda-\lambda^2=-u(1+u)
\]
is the exact idempotence defect of the non-unit Stone component.
\end{theorem}

\begin{proof}
The coefficient formula for $\mathcal Z_{1,u}$ is immediate from the definition. If one of $m,n$ equals $1$, then multiplicativity holds because $c_u(1)=1$. If $m,n>1$, then $mn>1$ and
\[
\Delta_u(m,n)=\lambda-\lambda^2.
\]
This proves the first formula.

For the Stone identity, note that
\[
\Delta(a_\lambda)=\Delta(\mathbf u)+\lambda\Delta(\mathbf v)
=\mathbf u\otimes\mathbf u+
\lambda(\mathbf v\otimes\mathbf u+\mathbf u\otimes\mathbf v+\mathbf v\otimes\mathbf v),
\]
while
\[
a_\lambda\otimes a_\lambda
=\mathbf u\otimes\mathbf u+
\lambda(\mathbf v\otimes\mathbf u+\mathbf u\otimes\mathbf v)+
\lambda^2\mathbf v\otimes\mathbf v.
\]
Subtracting yields the result. 
\end{proof}

\begin{theorem}\label{thm:coprime-kernel}
Define the interior coprime kernel by
\[
C(s,t):=\sum_{\substack{(m,n)=1\\ m,n>1}}m^{-s}n^{-t}.
\]
Then
\[
C(s,t)=\frac{\zeta(s)\zeta(t)}{\zeta(s+t)}-\zeta(s)-\zeta(t)+1.
\]
Consequently,
\[
\mathfrak D_u(s,t):=
\sum_{\substack{(m,n)=1}}\Delta_u(m,n)m^{-s}n^{-t}
=
-u(1+u)
\left(
\frac{\zeta(s)\zeta(t)}{\zeta(s+t)}-\zeta(s)-\zeta(t)+1
\right).
\]
In particular, at the doubled endpoint $u=1$,
\[
\mathfrak D_1(s,t)=
-2
\left(
\frac{\zeta(s)\zeta(t)}{\zeta(s+t)}-\zeta(s)-\zeta(t)+1
\right).
\]
\end{theorem}

\begin{proof}
The full coprime sum is
\[
\sum_{(m,n)=1}m^{-s}n^{-t}
=
\prod_p\left(1+\sum_{a\ge1}p^{-as}+\sum_{b\ge1}p^{-bt}\right)
=
\prod_p\frac{1-p^{-s-t}}{(1-p^{-s})(1-p^{-t})}
=
\frac{\zeta(s)\zeta(t)}{\zeta(s+t)}.
\]
Subtracting the cases $m=1$ or $n=1$ gives the formula for $C(s,t)$.

Now Theorem~\ref{thm:two-sheet-defect} shows that the coefficient-level defect equals $-u(1+u)$ exactly on the interior sector $m,n>1$ and vanishes otherwise. Hence
\[
\mathfrak D_u(s,t)=-u(1+u)C(s,t),
\]
which is the asserted formula. The case $u=1$ is immediate. 
\end{proof}

\begin{corollary}\label{cor:origin-values}
One has
\[
F_1(0)=-\frac32,
\qquad
\mathcal Z_{1,1}(0)=-2,
\qquad
\mathfrak D_1(0,0)=-3.
\]
Moreover,
\[
\mathfrak D_u(0,t)=u(1+u)F_1(t),
\qquad
\mathfrak D_1(0,t)=2F_1(t).
\]
\end{corollary}

\begin{proof}
Since $F_1(s)=\zeta(s)-1$,
\[
F_1(0)=\zeta(0)-1=-\frac32.
\]
Also,
\[
\mathcal Z_{1,1}(0)=2\zeta(0)-1=-2.
\]
Now
\[
C(0,0)=\frac{\zeta(0)\zeta(0)}{\zeta(0)}-\zeta(0)-\zeta(0)+1
=-\frac12+\frac12+\frac12+1
=\frac32,
\]
so
\[
\mathfrak D_1(0,0)=-2\cdot\frac32=-3.
\]
Further,
\[
C(0,t)=\frac{\zeta(0)\zeta(t)}{\zeta(t)}-\zeta(0)-\zeta(t)+1=1-\zeta(t),
\]
hence by Theorem~\ref{thm:coprime-kernel},
\[
\mathfrak D_u(0,t)=-u(1+u)(1-\zeta(t))=u(1+u)(\zeta(t)-1)=u(1+u)F_1(t).
\]
The case $u=1$ follows immediately. 
\end{proof}

\section{Number-field extension}

\begin{definition}\label{def:number-field-shift}
Let $K$ be a number field, let $a_K(n)$ be the number of nonzero ideals of $\cO_K$ of norm $n$, and define
\[
\zeta_K(s)=\sum_{n\ge1}a_K(n)n^{-s}.
\]
For $|t|<1$, define
\[
\zeta_{K,t}(s):=\sum_{n\ge1}a_K(n)(n+t)^{-s}
=
\sum_{0\neq J\triangleleft \cO_K}(N(J)+t)^{-s}.
\]
\end{definition}

\begin{theorem}\label{thm:number-field-deformation}
For $|t|<1$,
\[
\zeta_{K,t}(s)=\sum_{m\ge0}(-1)^m\frac{(s)_m}{m!}t^m\zeta_K(s+m),
\]
for $\Re(s)>1$, and this identity determines a meromorphic continuation of $\zeta_{K,t}$ to the whole plane. Moreover, $\zeta_{K,t}$ has the same unique simple pole at $s=1$ as $\zeta_K$, with the same residue.
\end{theorem}

\begin{proof}
Apply Theorem~\ref{thm:universal-taylor-deformation} and Corollary~\ref{cor:pole-rigidity} to the Dirichlet series $L=\zeta_K$, using the standard meromorphic continuation and pole structure of Dedekind zeta functions. 
\end{proof}

\begin{theorem}\label{thm:number-field-arithmetic-kernels}
For every integer $r\ge1$, the ratio
\[
Q_{K,r}(s):=\frac{\zeta_K(s+r)}{\zeta_K(s)}
\]
has an ideal-Dirichlet expansion
\[
Q_{K,r}(s)=\sum_{\mathfrak a\neq0}\alpha_{K,r}(\mathfrak a)N(\mathfrak a)^{-s},
\]
where $\alpha_{K,r}$ is multiplicative and
\[
\alpha_{K,r}(\mathfrak p^k)=-\frac{N(\mathfrak p)^r-1}{N(\mathfrak p)^{rk}}
\qquad(k\ge1).
\]
Equivalently,
\[
\boxed{
\alpha_{K,r}(\mathfrak a)=(-1)^{\omega(\mathfrak a)}
\frac{\prod_{\mathfrak p\mid \mathfrak a}(N(\mathfrak p)^r-1)}{N(\mathfrak a)^r}.
}
\]
\end{theorem}

\begin{proof}
Use the Euler product for $\zeta_K$:
\[
\zeta_K(s)=\prod_{\mathfrak p}(1-N(\mathfrak p)^{-s})^{-1}.
\]
Hence
\[
Q_{K,r}(s)=\prod_{\mathfrak p}\frac{1-N(\mathfrak p)^{-s}}{1-N(\mathfrak p)^{-s-r}}.
\]
The local computation is identical to the proof of Theorem~\ref{thm:arithmetic-kernels-riemann}, with $p$ replaced by $N(\mathfrak p)$. Multiplicativity follows from the Euler product. 
\end{proof}

\begin{corollary}\label{cor:number-field-non-eulerian}
If $0<t<1$, then $\zeta_{K,t}$ is not an ordinary Dirichlet series in the variable $n^{-s}$, hence has no classical Euler product.
\end{corollary}

\begin{proof}
The proof of Theorem~\ref{thm:non-eulerian-nonintegral-t} applies verbatim. The first nonzero coefficient of $\zeta_K$ occurs at $n=1$, since there is exactly one ideal of norm $1$, namely $\cO_K$. Thus the leading exponential scale of $\zeta_{K,t}(\sigma)$ is $(1+t)^{-\sigma}$, which cannot be the leading scale of an ordinary Dirichlet series in integer powers $n^{-\sigma}$. 
\end{proof}

\section{Final synthesis}

\begin{theorem}\label{thm:final-synthesis}
Relative to the baseline globalization note, the subsequent proved material changes the structure of $G(R)$ in the following exact ways.
\begin{enumerate}[label=(\arabic*)]
    \item \textbf{Universal quotients and characters.}
    The ring reflection $p_R:G(R)\to R$ is the universal ring quotient, and the support character $\chi_R:G(R)\to\B$ is the unique Boolean character.

    \item \textbf{Linear algebra.}
    The entire semimodule category is classified by support diagrams:
    \[
    G(R)\text{-}\SMod\simeq \Diag_R.
    \]
    Equivalently, an ordinary $R$-module is replaced by a join-semilattice with zero, an $R$-module on each support fiber, and transport maps along the order relation.

    \item \textbf{Cancellative sector.}
    The classical category survives exactly as the cancellative subcategory:
    \[
    G(R)\text{-}\SMod^{\mathrm{canc}}\simeq R\text{-}\Mod.
    \]

    \item \textbf{Free objects.}
    The free semimodule $G(R)^n$ is the Boolean cube of free $R$-modules:
    \[
    I\subseteq\{1,\dots,n\}
    \longmapsto
    R^I.
    \]

    \item \textbf{Products and multiplicative reconstruction.}
    Finite products are Boolean fiber products:
    \[
    G(A\times B)\cong G(A)\times_{\B}G(B).
    \]
    The multiplicative monoid alone does not reconstruct $G(R)$; the missing additive relations are exactly
    \[
    [a]+[b]=[a+b],
    \qquad
    [x]+[\tau]=[x].
    \]

    \item \textbf{Multiplicative pathology.}
    In the UFD case with infinitely many irreducibles, $\Spec(M_S)$ has infinite Krull dimension, and the contracted monoid algebra splits as
    \[
    \Z_0[M_S]\cong \Z\times \Z[R^{\bullet}].
    \]

    \item \textbf{Conservative ideal side.}
    Ideal multiplication remains classical:
    \[
    I_JI_K=I_{JK},
    \]
    and the ideal class monoid is unchanged. By contrast, the separated quotient-size weight $N(J)+1$ is not multiplicative, so the separated zeta sheet is non-Eulerian.

    \item \textbf{Element-theoretic factorization.}
    Units and supported irreducibles agree with those of $R$, but the supported zero $e=0_R$ becomes reducible and divisibility-prime.

    \item \textbf{All-orders shifted-zeta deformation.}
    The shifted quotient-size family is governed by the universal operator
    \[
    D=sT,
    \qquad
    (Tf)(s)=f(s+1),
    \]
    with formal flow
    \[
    L_t=e^{-tD}L,
    \qquad
    \partial_t^{\,n}L_t(s)=(-1)^n(s)_nL_t(s+n).
    \]
    Its logarithmic $n$th jet is a universal partition polynomial in the shifted ratios $L_t(s+m)/L_t(s)$.

    \item \textbf{Non-Eulerianity and endpoint defect.}
    For nonintegral $t$, the shifted family exits the Euler category altogether. At the endpoint $t=1$, the defect is controlled exactly by the two-point Stone scalar
    \[
    \delta(\lambda)=\lambda-\lambda^2.
    \]
    For the doubled sheet $2\zeta(s)-1$, one has
    \[
    \delta(2)=-2.
    \]

    \item \textbf{Arithmetic interpolation.}
    In the Hurwitz specialization,
    \[
    D_t(-m)=\frac{B_{m+1}(1+t)-B_{m+1}}{m+1},
    \qquad
    e^{-D_t'(0)}=\Gamma(1+t),
    \]
    so the deformation interpolates finite power sums and factorials.

    \item \textbf{Number fields.}
    All deformation-theoretic statements persist for Dedekind zeta functions, with the same pole rigidity and the same loss of Eulerianity for nonintegral deformation parameter.
\end{enumerate}
\end{theorem}

\begin{proof}
Items (1) through (12) are precisely Theorems~\ref{thm:ring-quotient}, \ref{thm:boolean-character-unique}, \ref{thm:full-faithfulness}, \ref{thm:equivalence-semimodules-diagrams}, \ref{cor:cancellative-category}, \ref{thm:free-semimodule-boolean-cube}, \ref{thm:product-fiber}, \ref{thm:monoid-reconstruction}, \ref{thm:monoid-prime-explosion}, \ref{thm:contracted-monoid-algebra}, \ref{thm:ideal-multiplication-classical}, Proposition~\ref{prop:separated-weight-nonmultiplicative}, Proposition~\ref{prop:units-irreducibles}, Theorem~\ref{thm:supported-zero-prime-reducible}, Theorem~\ref{thm:universal-taylor-deformation}, Corollary~\ref{cor:jets}, Theorem~\ref{thm:cumulant-partition-formula}, Theorem~\ref{thm:non-eulerian-nonintegral-t}, Theorems~\ref{thm:two-sheet-defect} and \ref{thm:coprime-kernel}, Corollaries~\ref{cor:defect-entire-bernoulli-gamma} and \ref{cor:faulhaber-factorial}, and Theorems~\ref{thm:number-field-deformation} and \ref{thm:number-field-arithmetic-kernels}. 
\end{proof}

\begin{thebibliography}{9}

\bibitem{DLMF}
NIST Digital Library of Mathematical Functions,
\emph{Chapter 25: Zeta and Related Functions}, especially \S25.11 on the Hurwitz zeta function,
\url{https://dlmf.nist.gov/25.11}.

\bibitem{MilneANT}
J.~S.~Milne,
\emph{Algebraic Number Theory}, available at
\url{https://www.jmilne.org/math/CourseNotes/ANT.pdf}.

\bibitem{Neukirch}
J.~Neukirch,
\emph{Algebraic Number Theory},
Grundlehren der Mathematischen Wissenschaften, vol.~322,
Springer, 1999.

\end{thebibliography}

\end{document}
